\documentclass[11pt,a4paper]{article}

% ─── Packages ─────────────────────────────────────────────────
\usepackage[T1]{fontenc}
\usepackage[utf8]{inputenc}
\usepackage{lmodern}
\usepackage[margin=2.2cm]{geometry}
\usepackage{xcolor}
\usepackage{titlesec}
\usepackage{enumitem}
\usepackage{fancyhdr}
\usepackage{parskip}
\usepackage{array}
\usepackage{booktabs}
\usepackage[hidelinks]{hyperref}

% ─── Colours ──────────────────────────────────────────────────
\definecolor{accent}{HTML}{1F6FEB}
\definecolor{ink}{HTML}{16181D}
\definecolor{soft}{HTML}{5B6270}
\definecolor{rule}{HTML}{D0D4DC}

% ─── Section styling ──────────────────────────────────────────
\titleformat{\section}
  {\large\bfseries\color{accent}}{\thesection.}{0.6em}{}
\titlespacing*{\section}{0pt}{1.3em}{0.5em}
\setlist[itemize]{leftmargin=1.3em,itemsep=2pt,topsep=3pt}
\setlist[enumerate]{leftmargin=1.5em,itemsep=3pt,topsep=3pt}

% ─── Header / footer ──────────────────────────────────────────
\pagestyle{fancy}
\fancyhf{}
\lhead{\small\color{soft}GovGuide \,$\cdot$\, Finals}
\rhead{\small\color{soft}AgenTriX 2026 \,$\cdot$\, Interim Submission}
\rfoot{\small\color{soft}\thepage}
\renewcommand{\headrulewidth}{0.4pt}
\renewcommand{\footrulewidth}{0pt}
\setlength{\headheight}{14pt}

% ─── Helpers ──────────────────────────────────────────────────
\newcommand{\lead}[1]{\textbf{\color{ink}#1}}
\newcommand{\subhead}[1]{\par\smallskip\noindent{\bfseries\color{accent}#1}\par\nopagebreak\smallskip}

\begin{document}
\thispagestyle{fancy}

% ─── Title block ──────────────────────────────────────────────
\begin{center}
  {\small\color{soft}\textsc{AgenTriX 2026 \,$\cdot$\, National-Level AI Hackathon \,$\cdot$\, Interim Submission (T+2)}}\\[0.7em]
  {\Huge\bfseries\color{ink}GovGuide}\\[0.35em]
  {\large\color{soft}An Agentic AI Navigator for Government \& Citizen Services}\\[0.9em]
  {\small\color{ink}Domain: Government \& Citizen Services \quad$\vert$\quad Finals \quad$\vert$\quad 20 June 2026}
\end{center}
\vspace{0.3em}
{\color{accent}\hrule height 1.6pt}
\vspace{0.4em}
{\color{rule}\hrule height 0.5pt}
\vspace{1.1em}

% ─── 1. Problem statement ─────────────────────────────────────
\section{Problem Statement}
Accessing a government service in Sri Lanka suffers from deep \lead{information asymmetry}. For any
service --- land deed transfer, NIC renewal, business registration, birth-certificate correction, or
a passport --- the exact documents to bring, the fees to pay, the correct office, the authorised
officer, and the processing sequence are known mainly to the staff at one counter on one day. These
requirements \lead{change over time, are applied inconsistently, and are almost never consolidated in one
current place online}; the official websites that exist are frequently incomplete, outdated, or hard
to navigate.

Consider a citizen transferring an \lead{inherited land deed} at their Divisional Secretariat. On the
first visit they are missing a certified extract; on the second, a copy of the death certificate; on
the third, a stamp-duty assessment issued at a different counter. Each failed visit is a lost day of
wages and another bus fare. This pattern --- discovering one more missing document each time ---
repeats across business registration, NIC renewals, vehicle revenue licences, and birth-certificate
corrections, each with its own overlapping departments and unclear sequencing.

As a result, a citizen often makes \lead{three or four separate failed visits} before succeeding. The
cost --- in time, money, and frustration --- is enormous, and it falls disproportionately on
\lead{rural citizens}, who absorb repeated transport costs for failed trips, and on those with
\lead{low literacy} in formal Sinhala or English, for whom even the available written guidance is
inaccessible. There is no single, current, plain-language, and \lead{personalised} source that tells a
citizen exactly what to do for \emph{their} specific case.

Three properties make this genuinely hard, and make naive solutions fail:
\begin{itemize}
  \item \lead{The domain is huge} --- hundreds of services across many departments; a knowledge base
        cannot be fully built up front, and it goes stale.
  \item \lead{Processes branch} --- e.g. land transfer differs for inheritance vs.\ sale vs.\ gift;
        a generic answer is wrong for the individual.
  \item \lead{Trust matters} --- a confidently wrong answer wastes a real trip, so answers must be
        grounded and attributable.
\end{itemize}

% ─── 2. Chosen domain ─────────────────────────────────────────
\section{Chosen Domain}
\lead{Government \& Citizen Services} (AgenTriX 2026 Domain~3), focused on Sri Lanka.

We chose this domain because the pain is universal, recurring, and measurable: nearly every adult
interacts with multiple government services over their lifetime, yet the process remains opaque even
to educated, persistent citizens. The domain spans a large and diverse set of high-frequency services,
including:
\begin{itemize}
  \item \lead{Civil registration} --- birth, marriage, and death certificates and corrections.
  \item \lead{Identity and travel} --- NIC issuance and renewal, passport applications.
  \item \lead{Property} --- land deed transfers, valuations, and registration.
  \item \lead{Business} --- business-name and company registration, and licences.
  \item \lead{Transport} --- vehicle revenue licences and registration.
\end{itemize}
Because these services share a common failure mode --- scattered and frequently
changing requirements --- a single intelligent layer that consolidates and personalises them can
relieve a disproportionate amount of citizen hardship, while remaining directly useful to the
government agencies that operate the counters.

% ─── 3. Solution outline ──────────────────────────────────────
\section{Solution Outline --- \textit{GovGuide}}
\lead{GovGuide} is an agentic AI citizen-services navigator. A citizen describes their need in plain
language, and the system works through the following steps:
\begin{enumerate}
  \item \lead{Identifies the correct service} from the natural-language request.
  \item Runs a short \lead{clarifying interview} --- asking only the questions needed to pin the exact
        sub-case, instead of giving a generic answer.
  \item \lead{Retrieves the current requirements} (documents, fees, the right office and officer, and
        an estimated timeline) from a curated knowledge base, with a \lead{cited source and a
        ``last-verified'' date} for every fact.
  \item If the knowledge base has a gap, autonomous agents \lead{research, curate, and write the answer
        back} to the knowledge base (our self-expanding RAG loop), then answer.
  \item Generates a \lead{personalised, printable Action Pack} --- the exact document checklist for the
        citizen's case, and an estimated cost breakdown.
\end{enumerate}

\subhead{System architecture at a glance}
GovGuide is built as a \lead{hierarchical multi-agent system}. A central Supervisor coordinates two
specialised teams: an \lead{Answering team} that handles the live conversation --- understanding the
request, identifying the service, conducting the clarifying interview, retrieving evidence, judging
whether that evidence is sufficient, and composing the answer --- and a \lead{Knowledge-Acquisition
team} that is invoked only when the knowledge base cannot answer. This separation keeps each agent
focused on a single responsibility and makes the system's behaviour transparent and auditable.

\subhead{Key innovation: a self-expanding knowledge base}
Because the government domain is far too large to index in advance, GovGuide does not attempt to
pre-load every service. Instead, its knowledge base is \lead{self-expanding}: when a citizen's query
hits a gap, the system autonomously researches the answer from trusted sources, curates it into
structured, attributable records, and writes it back to the knowledge base. The first citizen to ask
about a service effectively builds the knowledge for everyone who follows --- keeping the system both
fresh and continuously improving, and turning each query into a lasting asset.

\subhead{What the citizen receives}
The end product is not a wall of chat text but a concrete, take-away \lead{Action Pack}: a personalised
document checklist tailored to the citizen's exact case, together with an estimated cost breakdown, in
a clean printable format the citizen can carry to the counter. Every item is traceable to a cited,
dated source, so both the citizen and the officer can trust it.

\subhead{Example: an inherited land-deed transfer}
Returning to the citizen above: they type, \emph{``I want to transfer my late father's land to my
name.''} GovGuide identifies the service as a land deed transfer, then asks only the
few questions that matter --- such as the basis of transfer (inheritance) --- and nothing more. It
retrieves the inheritance-specific requirements; finding the stamp-duty step missing from the
knowledge base, it researches and fills that gap on the spot, then returns a single
Action Pack listing every document needed --- including the death certificate and certified extract
that previously caused failed visits --- each with an estimated cost and a citation. What once took
three or four trips becomes one informed visit.

% ─── 4. How we plan to apply AI ───────────────────────────────
\section{How We Plan to Apply AI}
AI is the core of GovGuide, applied across three layers --- orchestration, retrieval, and generation.

\subhead{Agentic orchestration with LangGraph}
The system is orchestrated as a \lead{stateful graph of cooperating agents} using LangGraph. A
Supervisor routes the conversation between the two teams, maintaining a single shared state so the
clarifying interview can pause for the citizen's input and resume seamlessly. Each agent owns one
responsibility:

\begin{center}
\small
\setlength{\tabcolsep}{6pt}
\renewcommand{\arraystretch}{1.18}
\begin{tabular}{@{}p{1.9cm} l p{8.0cm}@{}}
\toprule
\textbf{Team} & \textbf{Agent} & \textbf{Responsibility} \\
\midrule
Answering & A1 Intake \& Intent & Normalise the request; extract intent and key entities. \\
          & A2 Service Identifier & Map the request to a specific, known government service. \\
          & A3 Clarification & Ask the minimal questions to pin the exact variant of the service. \\
          & A4 Retrieval & Fetch the most relevant evidence from the knowledge base. \\
          & A5 Gap Grader & Judge whether the retrieved evidence is sufficient to answer. \\
          & A6 Action-Pack Gen. & Compose the cited, structured Action Pack for the citizen. \\
\midrule
Acquisition & B1 Research & Gather missing evidence from trusted sources on demand. \\
            & B2 Extract \& Curate & Structure raw findings into attributable records. \\
            & B3 KB Updater & Embed and write new knowledge back into the base. \\
            & B4 Moderation & Flag low-confidence knowledge for human review. \\
\bottomrule
\end{tabular}
\end{center}

\subhead{Self-expanding Retrieval-Augmented Generation (RAG)}
Retrieval is \lead{corrective}. After the Retrieval agent gathers candidate evidence, the Gap Grader
judges whether it actually answers the question. Only if it does not does the system trigger knowledge
acquisition --- searching a curated local source pool first and falling back to trusted public sources
--- within a \lead{bounded retry loop} that prevents runaway cost. Newly gathered knowledge is embedded
and indexed immediately, so retrieval steadily improves with use.

\subhead{Grounded, cited, and structured generation}
Every fact the system presents carries its source and a \lead{``last-verified'' date}, which both
builds citizen trust and curbs hallucination. The final Action Pack is produced as a
\lead{schema-validated structured object} rather than free text, so no required document is silently
dropped and the output renders reliably to a clean, printable checklist.

\subhead{Asking the fewest possible questions}
The clarifying interview is driven by the \lead{structure of each service's variants}: the system asks
only the questions whose answers actually change the requirements, and skips anything the citizen has
already implied. A citizen is therefore never burdened with irrelevant prompts, yet the system still
narrows down to a single, precise, case-specific answer --- the behaviour that turns a generic FAQ into
genuine, personalised guidance.

\subhead{Reliability by design}
Autonomy is deliberately \lead{bounded}. Knowledge acquisition runs inside a capped retry loop and
draws only from an allow-list of trusted sources, so the system cannot loop indefinitely or ingest
unreliable material. When it cannot confirm a fact, it says so plainly and labels the answer as
unverified rather than guessing. These guardrails are what make an AI that writes its own knowledge
base safe to place in front of citizens.

\subhead{Knowledge sources and trustworthiness}
The knowledge base is assembled from authoritative material --- gazette notifications, Divisional
Secretariat circulars, and official service-portal content --- supplemented by curated, real-world
input. Each record stores its origin and verification status, so newly acquired knowledge is clearly
distinguished from confirmed knowledge until a human reviewer promotes it. This provenance-first design
is what allows an autonomous, self-expanding system to remain trustworthy enough for government
information.

\subhead{Expected impact}
By converting scattered, counter-bound knowledge into a personalised, cited checklist, GovGuide aims to
collapse three or four failed visits into a single successful one --- saving citizens time, transport
cost, and lost wages, while easing counter congestion for the agencies. The benefit is largest
precisely for the rural and lower-literacy citizens who are least able to absorb the current cost.

\end{document}
